% ID: SEL-D0076-OA-E02
% ESTADO: cotejada. Segunda pasada visual y borrador compilado revisados.
% Fuente: Selectividad/sel_2008_matematicas_aplicadas/Mat CCSS 4 - Examen y criterios de correción.pdf
% SHA-256: f1040adb7b9fb01c26219800e6212b9a26b93bb9d5e933743d18ee6ac26e163b
% Página: 1; opción A; ejercicio original 2.

El beneficio de una empresa, en miles de euros, viene dado por la función
\[
B(x)=-3x^2+120x+675,\quad x\geq0
\]
donde $x$ representa el gasto en publicidad, en miles de euros.
\begin{enumerate}[label=\alph*)]
\item \textbf{(0.75 puntos)} Calcule el gasto a partir del cual la empresa no obtiene beneficios.
\item \textbf{(0.75 puntos)} Calcule el valor de $x$ que produce máximo beneficio. \textquestiondown{}Cuánto es ese beneficio?
\item \textbf{(0.75 puntos)} Determine los intervalos de crecimiento y decrecimiento del beneficio de la empresa.
\item \textbf{(0.75 puntos)} Represente gráficamente la función $B$.
\end{enumerate}
